\documentclass{article}
\usepackage[T1]{fontenc}
\usepackage{lmodern}
\usepackage{graphicx}
\usepackage{amsmath,amssymb}
\usepackage[a4paper,margin=2cm]{geometry}
\usepackage[absolute,overlay]{textpos}
\setlength{\TPHorizModule}{1mm}
\setlength{\TPVertModule}{1mm}
\usepackage[indonesian]{babel}
\usepackage{hyperref}
\setlength{\emergencystretch}{2em}
% unit: Halaman 34

\begin{document}

% === Halaman 34 (llm) ===

\section*{Cipher-Feedback (CFB)}

\begin{itemize}
    \item Mengatasi kekurangan pada mode $CBC$ apabila diterapkan pada pengiriman data yang belum mencapai ukuran satu blok
    \item Data dienkripsi dalam unit yang lebih kecil $(s)$ daripada ukuran blok $(b)$.
    \item Unit data yang dienkripsi panjangnya bisa 1 bit, 2 bit, 4-bit, 8 bit, dan lain-lain.
    \item Bila unit yang dienkripsi adalah satu karakter setiap kalinya, maka mode $CFB$-nya disebut $CFB$ 8-bit.
\end{itemize}

\end{document}
